\documentclass[10pt,a4paper]{amsart}
\usepackage[margin=19mm]{geometry}
\usepackage{amsmath,amssymb,mathtools}
\usepackage[scr=rsfs,cal=euler]{mathalfa}
\usepackage{xfrac}
\usepackage[unicode,hidelinks,bookmarks=false]{hyperref}
\newcommand{\ie}{i.e.\ }
\newcommand{\cf}{cf.\ }
\newcommand{\resp}{respectively\ }
\newcommand{\Subsection}{\subsection}
\providecommand{\nobreakdash}{\nobreak\hbox{-}}
\setlength{\emergencystretch}{2em}
\multlinegap0pt
\usepackage{amssymb}
\usepackage[shortlabels]{enumitem}
\usepackage[normalem]{ulem}
\newtheorem{cor}{Corollary}
\def\Ker{{\rm Ker}}
\def\Z{{\mathbb Z}}
\title{More CI-groups with respect to ternary relational structures}
\author{Ted Dobson, Joy Morris, Mikhail Muzychuk,, Pablo Spiga}
\date{Source: arXiv:2602.08802v1 (CC BY 4.0)}
\begin{document}
\maketitle

\noindent\emph{Source and licence.} More CI-groups with respect to ternary relational structures. Version arXiv:2602.08802v1 (CC BY 4.0), distributed by arXiv under the Creative Commons Attribution 4.0 International licence, http://creativecommons.org/licenses/by/4.0/, as recorded in the arXiv licence metadata for this exact version and shown on the abstract page https://arxiv.org/abs/2602.08802v1. Full licence notice: \texttt{licenses/arxiv-2602.08802-CC-BY-4.0.txt}.\par\smallskip

\noindent\textit{Editorial scope.} Counted unit: the article's cor cor2 (source0.tex lines 910--911). The article's complete original proof of it is delivered with it (source0.tex lines 913--948, one \texttt{proof} environment of 36 lines). No mathematics is written, repaired or re-derived: the statement and the proof are the article's own lines, in the article's own notation, and no retained line is substituted or re-flowed.  Retained uncounted as the source-local context the proof needs: lines 900--903.  Nothing was omitted from any retained span, so the package carries no omission record.  No retained line contains a citation, so the wrapper carries no bibliography and no bibliography entry.  No retained line contains a figure environment, an image-inclusion command or drawing code, so no image is delivered and the package declares no asset dependency.  Editorial formatting only, in addition: an AMS \texttt{amsart} preamble that restates the article's own declarations and macros used by the retained lines, namely the environments \texttt{cor} on separate counters, together with the standard packages those lines need.  The retained environments print in this document as Corollary 1 in order of appearance, whereas the article prints the same items with the numbering of its own full document.  The complete native source of the article as distributed in the arXiv e-print for this exact version is bundled unchanged as \texttt{sources/194275/source0.tex}.
	\begin{cor}\label{cor1} Let $p$ be a prime and let $\omega \in \Z_p^*$ be a primitive $n$th root of unity for some integer $n$ with $2 < n < p-1$. 
		Let $F = \Z_p \rtimes \langle \omega \rangle$ be the Frobenius group of order $np$. 
		Then the inner holomorph of $F$ is $3$-closed, and $F$ is not a CI-group with respect to ternary relational structures.
	\end{cor}

	\begin{cor}\label{cor2}Let $p$ be a prime,  let $n$ be a divisor of $p-1$ with $2 < n < p-1$, and let $\alpha\colon \Z_n\rightarrow \Z_p^*$ be a non-injective group homomorphism. Assume that $|\Ker(\alpha)|$ is even whenever $n$ is even.  Then the semidirect product $\Z_p\rtimes_\alpha\Z_n$ is not a CI-group with respect to ternary relational structures.
	\end{cor}

	\begin{proof}Let $\omega$ be a primitive $n$th root of unity in $\Z_p$.  
		Without loss of generality we may assume that 
		$\alpha\colon \langle \omega \rangle \rightarrow \langle \omega \rangle$ 
		is an endomorphism of $\langle \omega \rangle$, that is, 
		$\alpha(\omega^i)=\omega^{ai}$ for some $a\in \Z_n$.  
		It follows from the assumptions that $a\notin \Z_n^*$, and that $a$ is even whenever $n$ is even.  
		These conditions ensure the existence of $b\in \Z_n^*$ such that $a-b$ is invertible in $\Z_n$.
		
		Consider the Frobenius group 
		$F=\Z_p\rtimes \Z_n = \{(x,\omega^i)\,:\,i\in\Z_n,\ x\in\Z_p\}$.  
		By Corollary~\ref{cor1}, the group $F_L\cdot F_R\cong F\times F$ is $3$-closed.  
		Write elements of $F\times F$ as $[(x,\omega^i),(y,\omega^j)]$ with $x,y\in\Z_p$ and $i,j\in\Z_n$.  
		In this notation we have:
		$$
		\begin{array}{rcl}
			F_L & = & \{[(x,\omega^i),(0,1)] \mid x\in \Z_p,\ i\in\Z_n\},\\ 
			F_R & = & \{[(0,1),(y,\omega^j)] \mid y\in \Z_p,\ j\in\Z_n\},\\
			(F_L\cdot F_R)_e = \mathsf{Inn}(F) 
			& = & 
			\{[(x,\omega^i),(x,\omega^i)] \mid x\in \Z_p,\ i\in\Z_n\}.
		\end{array}
		$$
		It is easy to check that the subsets
		\begin{align*}
			G_1 &= \{[(x,\omega^{ai}),(0,\omega^{bi})] : x\in\Z_p,\ i\in\Z_n\},\\
			G_2 &= \{[(0,\omega^{bi}),(y,\omega^{ai})] : y\in\Z_p,\ i\in\Z_n\}
		\end{align*}
		are two regular subgroups of $F_L\cdot F_R$, each isomorphic to $G$.  
		Regularity follows from the fact that 
		\[
		G_1 \cap \mathsf{Inn}(F)
		= \{[(0,\omega^{ai}),(0,\omega^{bi})] : ai = bi,\ i\in\Z_n\}
		\]
		has size $1$, and likewise for $G_2$.  
		Moreover, $G_1$ and $G_2$ are not conjugate in $F_L\cdot F_R$ because their Sylow $p$-subgroups $(G_i)_p$, $i=1,2$, are normal in $F_L\cdot F_R$ and are distinct.
	\end{proof}

\end{document}
